\section{Question 1: Batch vs. Streaming Processing}
% Write your answer for Question 1 here.

When looking at the trade-offs between batch and streaming, the deciding factor is usually the \textbf{"half-life"} of the data's usefulness. Traditional batch systems, epitomized by \textbf{MapReduce}, are engineered for processing and generating large-scale datasets through a functional model that automatically handles parallelization, data distribution, and fault tolerance at scale[cite: 1, 12, 16]. This abstraction is ideal for high-throughput processing on bounded, static datasets; for instance, predicting inventory for a fresh food delivery service based on historical sales is highly efficient via an overnight MapReduce job, which can scale to process terabytes of data across thousands of commodity machines. However, low-latency processing becomes necessary when a system must react to \textbf{unbounded data} instantly to prevent operational failure. Frameworks like \textbf{Spark} improved upon MapReduce by utilizing in-memory \textit{Resilient Distributed Datasets (RDDs)}, enabling faster iterative processing and \textbf{micro-batching}. Yet, for truly dynamic operations—like rerouting delivery drivers during a sudden traffic jam—even a micro-batch delay of a few minutes renders the insight obsolete. The absolute necessity for \textbf{strict low-latency streaming} emerges in safety-critical contexts, such as \textbf{MRI-guided radiotherapy} (e.g., the TrackRAD challenge), where the system must track a moving tumor in real-time to adjust the radiation beam. In such scenarios, a fraction of a second of latency means irradiating healthy tissue instead of the target. Ultimately, moving beyond batch processing to a continuous streaming architecture introduces massive engineering complexity, especially regarding \textbf{state management}; therefore, it should be reserved for use cases where a delayed response is fundamentally an \textit{incorrect} response, while batch processing remains the more stable and cost-effective choice for all other applications.









